%File: LSE_capacity_appendices.tex
%
% Standalone appendices for LSE_capacity_na.tex.
% Equation numbers (36)--(53) are frozen via \tag{N} to match the joint compile.
% Cross-references to main-text equations are hard-coded as numerals.
%
\documentclass[letterpaper]{article} % DO NOT CHANGE THIS
\usepackage[submission]{aaai2026}  % DO NOT CHANGE THIS
\usepackage{times}  % DO NOT CHANGE THIS
\usepackage{helvet}  % DO NOT CHANGE THIS
\usepackage{courier}  % DO NOT CHANGE THIS
\usepackage[hyphens]{url}  % DO NOT CHANGE THIS
\usepackage{graphicx} % DO NOT CHANGE THIS
\urlstyle{rm} % DO NOT CHANGE THIS
\def\UrlFont{\rm}  % DO NOT CHANGE THIS
\usepackage{natbib}  % DO NOT CHANGE THIS AND DO NOT ADD ANY OPTIONS TO IT
\usepackage{caption} % DO NOT CHANGE THIS AND DO NOT ADD ANY OPTIONS TO IT
\frenchspacing  % DO NOT CHANGE THIS
\setlength{\pdfpagewidth}{8.5in} % DO NOT CHANGE THIS
\setlength{\pdfpageheight}{11in} % DO NOT CHANGE THIS

\usepackage{algorithm}
\usepackage{algorithmic}
\usepackage{amsmath}
\usepackage{amssymb}
\usepackage{bm}

\usepackage{newfloat}
\usepackage{listings}
\DeclareCaptionStyle{ruled}{labelfont=normalfont,labelsep=colon,strut=off}
\lstset{%
	basicstyle={\footnotesize\ttfamily},
	numbers=left,numberstyle=\footnotesize,xleftmargin=2em,
	aboveskip=0pt,belowskip=0pt,%
	showstringspaces=false,tabsize=2,breaklines=true}
\floatstyle{ruled}
\newfloat{listing}{tb}{lst}{}
\floatname{listing}{Listing}

\pdfinfo{
/TemplateVersion (2026.1)
}

\setcounter{secnumdepth}{1}

% Custom macros (same as main file)
\newcommand{\bxi}{\boldsymbol{\xi}}
\newcommand{\bx}{\bm{x}}
\newcommand{\avg}[1]{\langle #1 \rangle}
\newcommand{\e}{{\mathrm{e}}}
\newcommand{\dd}{{\mathrm{d}}}
\newcommand{\Bnet}{\beta}
\newcommand{\phieq}{\varphi_{\mathrm{eq}}}
\newcommand{\fret}{f_{\mathrm{ret}}}
\newcommand{\gmax}{g_{\max}}
\newcommand{\ac}{\alpha_c}
\newcommand{\acEd}{\alpha_c^{\mathrm{ed}}}
\newcommand{\p}{\partial}
\newcommand{\cas}{\textstyle\frac}

\title{\bfseries Supplementary Material:\\Appendices for ``Thermodynamic Approach and Saddle-Dominated Capacity of the LSE Dense Associative Memory''}

\author{Anonymous submission}
\affiliations{~}

\begin{document}

\maketitle

\appendix

\section{Appendix A: Derivation of the prefactor \\in the non-retrieval sum}

The Debye uniform asymptotic of the modified Bessel function~\citep[10.41.3]{olver2010nist}
gives the prefactor of Eq.~(13). Keeping the leading ($k = 0$) term of
the asymptotic series,
\begin{equation}
I_\nu(\nu z') \sim \frac{\e^{\nu\eta(z')}}{(2\pi\nu)^{1/2}(1+{z'}^2)^{1/4}},
\tag{36}
\end{equation}
where $z' = 2\Bnet N/(N-2)$ and
\begin{equation}
    \eta(z') = \sqrt{1+{z'}^2} + \ln\frac{z'}{1+\sqrt{1+{z'}^2}}\,.
    \tag{37}
    \label{eq:bessel_debye}
\end{equation}
$\eta(z')$ is expanded to $\mathcal{O}(N^{-1})$ because it is multiplied by
$\nu = N/2 - 1$ in the exponent:
\begin{equation}
    \eta(z') = \eta(2\Bnet) + \eta'(2\Bnet) \frac{4\Bnet}{N} + O(N^{-2})\,,
    \tag{38}
\end{equation}
where
\begin{align}
    \eta(2\Bnet)  &= 1 + 2\Bnet\varphi^* + \ln\varphi^*\,, \tag{39}\\
    \eta'(2\Bnet) &= \varphi^* + \frac{1}{2\Bnet}\,. \tag{40}
\end{align}
Then
\begin{equation}
    \nu \eta(z') = \tfrac{N}{2} \eta(2\Bnet) - \ln\varphi^* + \mathcal{O}(1/N)\,.
    \tag{41}
\end{equation}
The $\mathcal{O}(1)$ piece $-\ln\varphi^*$ in the exponent contributes a factor
$1/\varphi^*$ to the prefactor. Using Stirling's formula
\begin{equation}
    \Gamma(N/2) \sim \sqrt{4\pi/N}\,(N/2)^{N/2}\,\e^{-N/2}
    \tag{42}
\end{equation}
and collecting all terms in Eq.~(18) to $\mathcal{O}(1)$ in the
exponent and to leading order in the prefactor reproduces the exponent
of Eq.~(13) and gives the prefactor:
\begin{equation}
S_{\rm nr} \approx \frac{\e^{N[\alpha + g_{\max}(\Bnet) - \Bnet]}}{\sqrt{1-\varphi^{*4}}}\,,
\tag{43}
\end{equation}
with $g_{\max}(\Bnet)$ given by Eq.~(12).
At $\Bnet = 1$, $\varphi^* \approx 0.618$ and the prefactor is $\approx 1.082$, reproducing Eq.~(14).


\section{Appendix B: Tangency of the cusp lines \\to the spherical extreme}

To show that the cusp lines~(16) are tangent to the
spherical-extreme curve~(6) at every
$\Bnet$, we need to show that the slope and the $\varphi$-intercept
coincide.

Differentiating~(6) once gives
\begin{equation}
\varphi_{\max}'(\alpha) \;=\; \frac{1-\varphi_{\max}^2}{\varphi_{\max}}\,,
\tag{44}
\label{eq:tangent_slope_C}
\end{equation}
and a second differentiation gives $\varphi_{\max}''(\alpha) < 0$ for $\alpha > 0$,
so $\varphi_{\max}$ is strictly concave on $(0, \infty)$. For a concave function
every tangent line lies above the function, with equality only at the tangency
point. This is used in the lower-envelope identity below.

The cusp line at sharpness $\Bnet$,
\begin{equation}
\varphi \;=\; \alpha/\Bnet + g_{\max}(\Bnet)/\Bnet\,,
\tag{45}
\end{equation}
has slope $1/\Bnet$ and $\varphi$-intercept $g_{\max}(\Bnet)/\Bnet$. To find the point on
the extreme curve where the tangent has the same slope, set the right-hand side
of~\eqref{eq:tangent_slope_C} equal to $1/\Bnet$:
\begin{equation}
\frac{1-\varphi^2}{\varphi} \;=\; \frac{1}{\Bnet}
\quad\Longleftrightarrow\quad
\Bnet\,\varphi^2 + \varphi - \Bnet = 0\,.
\tag{46}
\label{eq:saddle_eq_C}
\end{equation}
This is the saddle equation and its unique positive root is the spherical saddle~(11).
The corresponding $\alpha$ on the extreme curve follows from
$\varphi^* = \varphi_{\max}(\alpha^*)$:
\begin{equation}
\alpha^*(\Bnet) \;=\; -\tfrac12 \ln(1-\varphi^{*\,2})\,.
\tag{47}
\label{eq:alpha_star_C}
\end{equation}

A line of slope $1/\Bnet$ passing through $(\alpha^*(\Bnet), \varphi^*(\Bnet))$ has
$\varphi$-intercept
\begin{equation}
b_{\rm tg}(\Bnet) \;=\; \varphi^* - \alpha^*/\Bnet\,.
\tag{48}
\end{equation}
The cusp line has $\varphi$-intercept $g_{\max}(\Bnet)/\Bnet$.
Using~(12), the $\varphi$-intercepts coincide;
the two lines are therefore identical.

By the concavity statement above, the cusp line at any sharpness $\Bnet$ lies above the extreme curve, with equality only at
$\alpha = \alpha^*(\Bnet)$:
\begin{equation}
\varphi_c(\alpha, \Bnet) \;\ge\; \varphi_{\max}(\alpha) \quad \text{for all $\alpha$ and all $\Bnet > 0$}\,,
\tag{49}
\end{equation}
with equality at one $\alpha$ for each $\Bnet$. Taking the infimum over $\Bnet$,
\begin{equation}
\varphi_{\max}(\alpha) \;=\; \inf_{\Bnet > 0} \varphi_c(\alpha, \Bnet)\,,
\tag{50}
\label{eq:lower_envelope_C}
\end{equation}
the lower-envelope identity.

The saddle function $\varphi^*(\Bnet)$~(11) is continuous and strictly increasing on $(0, \infty)$ from $0$ to $1$. As $\Bnet$ sweeps $(0, \infty)$,
the tangency point $(\alpha^*(\Bnet), \varphi^*(\Bnet))$ therefore traces the entire extreme curve $\varphi_{\max}(\alpha)$, from $(0, 0)$ at $\Bnet \to 0$ to $\varphi_{\max} \to 1$ at $\Bnet \to \infty$.

Two results follow from the geometric setup above.  At $T = 0$ the
cusp line of slope $1/\Bnet$ reaches $\varphi = 1$ at
$\alpha = \Bnet - g_{\max}(\Bnet)$, which is the zero-temperature
capacity~(17) (equivalently, the spinodal at
$T = 0$).

The condition of \citet{ramsauer2021hopfield},
$\varphi_{\max}(\alpha) = 1 - \alpha/\Bnet$, is the intersection of
$\varphi_{\max}$ with the line of slope $-1/\Bnet$ through $(0, 1)$.
By~\eqref{eq:lower_envelope_C}, $\varphi_{\max}$ lies below every
cusp line, so this line crosses $\varphi_{\max}$ at strictly
smaller $\alpha$ than the cusp line at slope $1/\Bnet$ reaches
$\varphi = 1$.  Their capacity $\alpha_c^R$ therefore satisfies
$\alpha_c^R < \alpha_c(\Bnet, 0)$ for every $\Bnet > 0$.

Using~(11) in
(12) and then in~(17), one obtains:
\begin{equation}
\begin{aligned}
\ac(\Bnet, 0) = \Bnet - g_{\max}(\Bnet) &\;\to\; \Bnet
   && (\Bnet \to 0)\,,\\
\ac(\Bnet, 0) = \Bnet - g_{\max}(\Bnet) &\;\to\; \tfrac{1}{2}\ln\Bnet + \tfrac{1}{2}
   && (\Bnet \to \infty)\,,
\end{aligned}
\tag{51}
\label{eq:asymptotic_C}
\end{equation}
the small- and large-$\Bnet$ limits of the zero-temperature capacity
cited in the main text.

\paragraph{Legendre transform.}
The lower-envelope identity~\eqref{eq:lower_envelope_C} and the formula
below are the two sides of a Legendre transform between
$\varphi_{\max}(\alpha)$ and $g_{\max}(\Bnet)/\Bnet$.  The forward
transform takes $\varphi_{\max}$ to its conjugate,
\begin{equation}
g_{\max}(\Bnet)/\Bnet \;=\; \sup_\alpha\,\bigl[\,\varphi_{\max}(\alpha) - \alpha/\Bnet\,\bigr]\,,
\tag{52}
\label{eq:legendre_C}
\end{equation}
with the supremum attained at $\alpha = \alpha^*(\Bnet)$ from the
slope-match condition~\eqref{eq:saddle_eq_C}.  The Legendre transform
is involutive: applying it to $g_{\max}(\Bnet)/\Bnet$ recovers
$\varphi_{\max}(\alpha)$, which is exactly~\eqref{eq:lower_envelope_C}
after using the cusp-line definition~(16).

Differentiating~(12) in $\Bnet$ and using the
envelope theorem $\p g/\p\varphi|_{\varphi^*} = 0$ gives
\begin{equation}
g_{\max}'(\Bnet) \;=\; \varphi^*(\Bnet)\,,
\tag{53}
\label{eq:envelope_C}
\end{equation}
so $g_{\max}(\Bnet)$ is recovered by integrating $\varphi^*(\Bnet)$.
Plugging the small- and large-$\Bnet$ limits $\varphi^* \sim \Bnet$
and $\varphi^* \sim 1 - 1/(2\Bnet)$ into~\eqref{eq:envelope_C} and
integrating reproduces~\eqref{eq:asymptotic_C}.


\section{Appendix C: Translation of \citeauthor{ramsauer2021hopfield}'s storage-capacity bound}

We give a literal symbol-by-symbol map between
\citeauthor{ramsauer2021hopfield}'s notation and ours, then derive
the Theorem~4 reading, Eq.~(19) and
Eq.~(20) step by step.

\paragraph{Notation map.}
\citeauthor{ramsauer2021hopfield} use $N$ patterns
$\bm x_i \in \mathbb R^d$ on a sphere of radius
$\mathbf M := K\sqrt{d-1}$ with sharpness $\beta$, and a query
$\bm \xi \in \mathbb R^d$.  Our patterns are
$\bxi^\mu \in \mathbb R^N$, $\mu = 1, \dots, M$, on a sphere of
radius $\sqrt N$ with sharpness $\Bnet$; the target is $\bxi^1$ and
the query is taken at the target.  Setting $K = 1$ in
\citeauthor{ramsauer2021hopfield}'s convention identifies the
pattern norm $\mathbf M = \sqrt{d-1}$ with our $\sqrt N$ in the
large-$N$ limit.  The variable map is
\begin{center}
\begin{tabular}{lll}
\hline
\citeauthor{ramsauer2021hopfield} & here & meaning \\
\hline
$N$ & $M = \e^{\alpha N}$ & number of patterns \\
$d$ & $N$ & pattern dimension \\
$\mathbf M$ & $\sqrt N$ & pattern norm \\
$\beta$ & $\Bnet$ & softmax sharpness \\
$\bm x_i^\top \bm x_j$ & $N\,\varphi^{ij}$ & pair product \\
$\bm x_i$, $\bm \xi$ & $\bxi^i$, $\bxi^1$ & patterns, target query \\
\hline
\end{tabular}
\end{center}
With the query at the target, $\bm \xi^\top \bm x_1 = N$ and
$\bm \xi^\top \bm x_j = N\,\varphi^{1j}$ for $j \geq 2$.  The
\emph{maximum competing alignment} is
$\varphi_{\max} = \max_{j \neq 1}\varphi^{1j}$.

\paragraph{Theorem~4 reading: retrieval ties to $N(1-\varphi_{\max})$.}
The separation, their Eq.~(5), is
\begin{equation*}
\Delta_i \;:=\; \bm x_i^\top \bm x_i - \max_{j \neq i}\bm x_i^\top \bm x_j.
\end{equation*}
At the target $i = 1$ we substitute the map: $\bm x_1^\top \bm x_1 = N$
and $\max_{j \neq 1}\bm x_1^\top \bm x_j = N\,\varphi_{\max}$.  Hence
\begin{equation*}
\Delta_1 \;=\; N\,(1 - \varphi_{\max}).
\end{equation*}
Their Theorem~5, Eq.~(9), bounds the retrieval error of pattern
$\bm x_1$ by
\begin{equation*}
\|\bm x_1 - \bm x_1^*\| \;\leq\; 2 e\,(N-1)\,\mathbf M\,\e^{-\beta \Delta_1}.
\end{equation*}
Substituting the map ($N \to M$, $\mathbf M \to \sqrt N$,
$\beta \to \Bnet$, $\Delta_1 \to N(1-\varphi_{\max})$):
\begin{equation*}
\|\bxi^1 - \bxi^{1*}\| \;\leq\; 2 e\,M\,\sqrt N\,\e^{-\Bnet N(1-\varphi_{\max})}.
\end{equation*}
For this bound to guarantee retrieval (error below any fixed
$\epsilon$ at large $N$) we need the exponential to overcome the
prefactor, i.e.
\begin{equation*}
\Bnet N (1-\varphi_{\max}) \;>\; \ln M + O(\ln N) \;=\; \alpha N + O(\ln N),
\end{equation*}
since $M = \e^{\alpha N}$.  To leading order in $N$ this is
$\varphi_{\max} < 1 - \alpha/\Bnet$.  As $\varphi_{\max} \to 1$ the
exponential factor approaches $1$, the bound no longer guarantees
small retrieval error, and Theorem~4 stops securing retrieval.  In
this sense Theorem~4 ties retrieval to the separation
$N(1-\varphi_{\max})$, and retrieval fails when
$\varphi_{\max} \to 1$.

\paragraph{Derivation of Eq.~(19).}
The non-retrieval sum~(8) is
\begin{equation*}
S_{\rm nr} \;=\; \sum_{\mu = 2}^M \e^{\Bnet N (\varphi^\mu - 1)}.
\end{equation*}
For every $\mu \geq 2$, $\varphi^\mu \leq \varphi_{\max}$ by
definition of $\varphi_{\max}$, so
\begin{equation*}
\e^{\Bnet N (\varphi^\mu - 1)} \;\leq\; \e^{\Bnet N (\varphi_{\max} - 1)}
   \;=\; \e^{-\Bnet N (1-\varphi_{\max})}.
\end{equation*}
The sum of $M - 1$ terms, each at most this maximum, satisfies
\begin{equation*}
S_{\rm nr} \;\leq\; (M-1)\,\e^{-\Bnet N (1-\varphi_{\max})}
   \;\leq\; M\,\e^{-\Bnet N (1-\varphi_{\max})}.
\end{equation*}
Using $M = \e^{\alpha N}$ on the right,
\begin{equation*}
M\,\e^{-\Bnet N (1-\varphi_{\max})}
   \;=\; \e^{\alpha N - \Bnet N (1-\varphi_{\max})}
   \;=\; \e^{N(\alpha - \Bnet(1-\varphi_{\max}))},
\end{equation*}
which is Eq.~(19).  The $M$-times-largest step
is the same one \citeauthor{ramsauer2021hopfield} use in their
Lemma~A4 Eq.~(134) for the softmax probability mass on non-retrieval patterns;
applying it to the unnormalised shifted sum gives the bound on
$S_{\rm nr}$ above.

\paragraph{Derivation of Eq.~(20).}
The shifted softmax sum~(7) reads
$S(\bm x) = 1 + S_{\rm nr}$ at the target, since the retrieval
term equals $\e^{\Bnet N(\varphi^1 - 1)} = 1$ for $\varphi^1 = 1$.
Retrieval requires the retrieval term to dominate, $S_{\rm nr} < 1$;
substituting the bound just derived,
\begin{equation*}
S_{\rm nr} < 1 \ \ \Leftarrow\ \ 
\e^{N(\alpha - \Bnet(1-\varphi_{\max}))} < 1
\ \ \Leftrightarrow\ \ 
\alpha < \Bnet\,(1-\varphi_{\max}),
\end{equation*}
or $\varphi_{\max} < 1 - \alpha/\Bnet$, the same threshold as in the
Theorem~4 reading.  The capacity is where this inequality saturates,
$\varphi_{\max}(\alpha) = 1 - \alpha/\Bnet$.  Substituting the
exact-density extreme alignment~(6),
$\varphi_{\max}(\alpha) = \sqrt{1 - \e^{-2\alpha}}$,
\begin{equation*}
\sqrt{1 - \e^{-2\alpha}} \;=\; 1 - \frac{\alpha}{\Bnet},
\end{equation*}
which is Eq.~(20).  The Gaussian alternative
$\varphi_{\max}(\alpha) = \sqrt{2\alpha}$, used in
\citet{icml2026theory}, instead gives
$\sqrt{2\alpha} = 1 - \alpha/\Bnet$; at $\Bnet \to \infty$ this
reduces to $\ac = 1/2$.



\bibliography{aaai2026}

\end{document}
